%=============================================================
% INTRODUCTION
%=============================================================

A microcontroller that runs a convolutional network usually has no fan,
draws on a coin cell or a harvested supply, and shares its single core
with the other tasks of a real-time operating system
(RTOS)~\cite{tsoukas2024tinyml,ray2022tinyml}. Energy per inference limits
such a device more than anything else, and two groups of decisions set it.
Depth, channel width, and weight precision fix how many
multiply--accumulate (MAC) operations the network needs and how much memory
it occupies. The time slice, the priority policy, and the clock frequency
then fix how quickly, and at what power, the RTOS carries those operations
out. The groups interact: a lower clock saves power but stretches every
inference, and a smaller network frees latency headroom that a lower clock
can spend.

Tool flows usually tune one group while holding the other fixed. Neural
architecture search~\cite{ren2022nas}, in its hardware-aware form,
co-explores network topology against a given accelerator, compiler, or
device~\cite{chitty2023hwnas,jiang2020coexploration}, and on
microcontrollers it can treat memory and latency as constraints~\cite{king2025mcu}.
MCUNet co-designs the network with its inference library but leaves the RTOS
scheduler and the core clock out of the search~\cite{lin2023tinyml}.
Co-search of the network with a compiler-level scheduling policy
exists~\cite{chen2023chanas}; the RTOS parameters of the core, which are
time slice, priority, and clock, are outside it. Energy-aware scheduling has
been studied for FreeRTOS-based systems~\cite{oliveira2023freertos}, and
learned frequency scaling has been applied to inference
workloads~\cite{zhang2024dvfo}; there, the network is an input to the problem
and is never searched.

Bayesian optimization (BO) fits problems where each evaluation is costly and
no gradient is available~\cite{wang2023bo}. Multi-objective BO has been
applied to tinyML network search~\cite{deutel2025mobo}, and hypervolume-based
acquisition functions have been run with constraints and on discrete design
spaces~\cite{low2024constrained,mamun2025mobo}. None of the studies cited here
puts the scheduler parameters in the same search space as the architecture,
or tests whether searching them together beats searching them in turn.

This study does both. Network architecture, weight precision, and scheduler
configuration form one six-variable problem, accuracy against energy per
inference under accuracy, latency, deadline-miss, and memory constraints.
Accuracy comes from trained and quantized networks, cycles and memory from a
model calibrated on an STM32 board, and scheduling from a deterministic
simulation of a fixed-priority RTOS. Full enumeration of $12{,}000$ designs
supplies the reference front, and six search methods are scored against it
over 20 seeds. The contributions are the joint formulation with its
calibrated evaluation, a comparison that includes a two-stage ablation in
which the scheduler is fixed before the network is searched, and a
held-out check of the cost models on eight networks that were not used to
calibrate them.
